\section{Versuche}
\newpage

\subsection{Versuch 1: Elektrochemische Spannungsreihe}
\subsubsection{Aufgabenstellung}
Magnesium, Zink, Nickel und Kupfer sollen mittels ihrer Reaktion mit Salzsäure in die elektrochemische Spannungsreihe eingeordnet werden.
\subsubsection{Durchführung}


\begin{enumerate}[label=\textbf{\alph*)}]
    \item Vier Reagenzgläser wurden mit einer halbkonzentrierter Salzsäure befüllt. In diese wurden anschließend jeweils etwas Magnesiumband, Zinkgranalie, Nickelpulver und Kupferspäne gegeben.
    \item Drei Reagenzgläser wurden mit Zinkgranalie befüllt, drei weitere mit Kupferspänen. Die beiden Metalle wurde nun paarweise konzentrierte Salzsäure, verdünnte Salpetersäure und konzentrierte Salpetersäure versetzt.
\end{enumerate}

\subsubsection{Beobachtung}



\begin{enumerate}[label=\textbf{\alph*)}]
    \item
    \item
\end{enumerate}

\subsubsection{Auswertung}
\subsubsection{Fehlerbetrachtung}

\subsection{Versuch 2: Abhängigkeit des Redoxpotentials vom pH-Wert}
\subsubsection{Aufgabenstellung}
\subsubsection{Durchführung}
\subsubsection{Beobachtung}
\subsubsection{Auswertung}
\subsubsection{Fehlerbetrachtung}

\subsection{Versuch 3: Redox-Amphoterie, Dis- und Synproportionierung}
\subsubsection{Aufgabenstellung}
\subsubsection{Durchführung}
\subsubsection{Beobachtung}
\subsubsection{Auswertung}
\subsubsection{Fehlerbetrachtung}

\subsection{Versuch 4: Redox-Titration einer $\ce{Cu^{2+}}$-Kationen-Lösung}
\subsubsection{Aufgabenstellung}
\subsubsection{Durchführung}
\subsubsection{Beobachtung}
\subsubsection{Auswertung}
\subsubsection{Fehlerbetrachtung}




\section{Zusammenfassung}
